\documentclass[10pt]{article}
\providecommand{\FIGDIR}{fig}
\newcommand{\DIGESTDATE}{18 August 2026}
% ===================== PM Research Watch - shared preamble =====================
\usepackage[T1]{fontenc}
\usepackage[utf8]{inputenc}
\usepackage{amsmath,amssymb}
\usepackage{tgpagella}
\usepackage{tgheros}
\renewcommand{\sfdefault}{qhv}
\renewcommand{\ttdefault}{lmtt}
\usepackage{microtype}
\usepackage[english]{babel}

\usepackage[a4paper,top=17mm,bottom=17mm,left=17mm,right=17mm,
            headsep=6mm,footskip=10mm]{geometry}
\usepackage{graphicx}
\usepackage{xcolor}
\usepackage{colortbl}
\usepackage{tikz}
\usetikzlibrary{calc}
\usepackage[most]{tcolorbox}
\usepackage{titlesec}
\usepackage{fancyhdr}
\usepackage{booktabs}
\usepackage{array}
\usepackage{longtable}
\usepackage{enumitem}
\usepackage{multicol}
\usepackage{ragged2e}
\usepackage[hidelinks]{hyperref}

% ------------------------------------------------------------------ palette
\definecolor{Ink}    {HTML}{14262E}
\definecolor{Deep}   {HTML}{0F4C5C}
\definecolor{Teal}   {HTML}{1F7A8C}
\definecolor{Sky}    {HTML}{4FA3C4}
\definecolor{Sage}   {HTML}{6FA287}
\definecolor{Amber}  {HTML}{E3A008}
\definecolor{Clay}   {HTML}{B5643C}
\definecolor{Coral}  {HTML}{D1495B}
\definecolor{Violet} {HTML}{6A5B9E}
\definecolor{Slate}  {HTML}{4A6572}
\definecolor{Paper}  {HTML}{FBF9F5}
\definecolor{Mist}   {HTML}{EEEAE2}
\definecolor{Rule}   {HTML}{D9D3C9}

\pagecolor{Paper}
\color{Ink}

\hypersetup{colorlinks=true, linkcolor=Deep, urlcolor=Teal, citecolor=Teal,
            pdftitle={PM Research Watch}, pdfauthor={Automated daily digest}}

% ------------------------------------------------------------------ headings
\titleformat{\section}
  {\sffamily\bfseries\large\color{Deep}}{}{0pt}
  {}[{\vspace{1.2mm}\color{Amber}\titlerule[1.1pt]}]
\titlespacing*{\section}{0pt}{5.5mm}{2.8mm}

\titleformat{\subsection}
  {\sffamily\bfseries\normalsize\color{Teal}}{}{0pt}{}
\titlespacing*{\subsection}{0pt}{3.6mm}{1.4mm}

\setlist[itemize]{leftmargin=4.2mm, itemsep=1.1mm, topsep=1.3mm, parsep=0pt,
                  label=\textcolor{Teal}{\footnotesize$\blacktriangleright$}}

% ------------------------------------------------------------------ header/footer
\pagestyle{fancy}
\fancyhf{}
\renewcommand{\headrulewidth}{0pt}
\renewcommand{\footrulewidth}{0pt}
\fancyhead[L]{\sffamily\scriptsize\color{Slate}PM RESEARCH WATCH\;\textcolor{Rule}{\textbar}\;\DIGESTDATE}
\fancyhead[R]{\sffamily\scriptsize\color{Slate}Particulate matter monitoring \& health effects}
\fancyfoot[C]{\sffamily\scriptsize\color{Slate}\thepage}

% ------------------------------------------------------------------ components
% metric tile
\newcommand{\metric}[3]{%
  \begin{tcolorbox}[enhanced, colback=white, colframe=Rule, boxrule=0.5pt,
    arc=1.6mm, left=2.2mm, right=2.2mm, top=1.8mm, bottom=1.8mm,
    borderline west={1.7mm}{0pt}{#3}, nobeforeafter, valign=center,
    height=17mm, width=\linewidth]
    \RaggedRight
    {\sffamily\bfseries\fontsize{17}{17}\selectfont\color{#3}#1}\\[0.5mm]
    {\sffamily\fontsize{6.5}{7.4}\selectfont\color{Slate}#2}
  \end{tcolorbox}}

% pull-quote / key insight box
\newtcolorbox{keybox}[1][]{enhanced, colback=Mist, colframe=Mist,
  arc=1.6mm, left=3.2mm, right=3.2mm, top=2.4mm, bottom=2.4mm,
  borderline west={1.7mm}{0pt}{Amber}, #1}

% subtopic band
\newcommand{\band}[2]{%
  \vspace{3.2mm}\noindent
  \begin{tcolorbox}[enhanced, colback=#2, colframe=#2, arc=1.2mm,
    left=2.6mm, right=2.6mm, top=1.1mm, bottom=1.1mm, boxrule=0pt,
    width=\linewidth, nobeforeafter]
    {\sffamily\bfseries\small\color{white}#1}
  \end{tcolorbox}\vspace{1.6mm}}

% one paper entry: {short cite}{title}{journal}{doi}{body}
\newcommand{\paperentry}[5]{%
  \par\vspace{1.9mm}\noindent
  {\sffamily\bfseries\footnotesize\color{Ink}#1}\;%
  {\sffamily\fontsize{7}{8}\selectfont\color{Slate}#3}\\[0.4mm]
  {\footnotesize\itshape\color{Deep}#2}\\[0.6mm]
  {\footnotesize #4}\\[0.3mm]
  {\fontsize{6.6}{7.6}\selectfont\color{Slate}\texttt{doi:}\,\href{https://doi.org/#5}{\texttt{\detokenize{#5}}}}
  \par}

% caption style
\newcommand{\figcap}[1]{\par\vspace{1.0mm}
  {\fontsize{7.2}{8.6}\selectfont\color{Slate}\textsf{\textbf{Fig.}}\; #1\par}\vspace{1.5mm}}

\providecommand{\permil}{\textperthousand}

\setlength{\parindent}{0pt}
\setlength{\parskip}{0pt}
\renewcommand{\arraystretch}{1.18}


\begin{document}

% ==================================================================== MASTHEAD
\thispagestyle{empty}
\begin{tikzpicture}[remember picture, overlay]
  \fill[Deep] (current page.north west) rectangle
      ($(current page.north east)+(0,-67mm)$);
  \fill[Amber] ($(current page.north west)+(0,-67mm)$) rectangle
      ($(current page.north east)+(0,-68.4mm)$);
  \foreach \x/\y/\r in {22/12/0.9, 41/26/1.5, 63/9/0.7, 88/31/1.1, 112/17/1.9,
                        138/29/0.8, 159/13/1.3, 178/54/1.0, 196/21/1.6,
                        30/44/0.6, 74/57/1.2, 125/48/0.9, 168/6/0.8, 191/59/0.7,
                        52/18/1.0, 100/60/0.6, 148/20/1.1, 66/49/1.4, 205/45/0.8}
    \fill[Sky, opacity=0.30] ($(current page.north west)+(\x mm,-\y mm)$) circle (\r mm);
\end{tikzpicture}

\vspace*{4mm}
{\sffamily\fontsize{7.6}{9}\selectfont\color{Sky}%
 AUTOMATED DAILY LITERATURE DIGEST\;\textbullet\;PhD RESEARCH WATCH\par}
\vspace{2.2mm}
{\sffamily\bfseries\fontsize{27}{29}\selectfont\color{white}%
 Particulate Matter\par}
{\sffamily\fontsize{27}{29}\selectfont\color{Sky}%
 Monitoring \& Health Effects\par}
\vspace{5.0mm}
{\sffamily\fontsize{8.4}{10}\selectfont\color{white}%
 Issue 2026-08-18\;\;\textcolor{Amber}{\textbar}\;\;Entry window 18 August 2026\;\;%
 \textcolor{Amber}{\textbar}\;\;18 records carried, 12 rejected\par}
\vspace{18mm}

% -------------------------------------------------------------- metric strip
\noindent
\begin{minipage}[t]{0.190\linewidth}\metric{18}{RECORDS CARRIED\\(12 REJECTED)}{Deep}\end{minipage}\hfill
\begin{minipage}[t]{0.190\linewidth}\metric{8}{SUBTOPIC CLUSTERS\\REPRESENTED}{Teal}\end{minipage}\hfill
\begin{minipage}[t]{0.190\linewidth}\metric{9}{EFFECT ESTIMATES\\WITH INTERVALS}{Sky}\end{minipage}\hfill
\begin{minipage}[t]{0.190\linewidth}\metric{9}{INSTRUMENTATION /\\EXPOSURE RECORDS}{Amber}\end{minipage}\hfill
\begin{minipage}[t]{0.190\linewidth}\metric{2}{METADATA-ONLY\\(LABELLED)}{Clay}\end{minipage}

\vspace{6mm}

% -------------------------------------------------------------- signal of day
\section*{\sffamily\bfseries\color{Deep}Signal of the day}
\vspace{-1.5mm}{\color{Amber}\rule{\linewidth}{1.1pt}}\vspace{2.5mm}

\begin{keybox}
{\sffamily\bfseries\small\color{Deep}A single \textit{ES\&T} issue supplies both halves of
the calibration problem this watch tracks --- the network-scale fix and the reason the
network's own target quantity is the wrong one.}\\[1.4mm]
\footnotesize
Wang et al.\ remove the colocation requirement from low-cost sensor calibration. Their
framework pairs uncolocated sensors to reference stations by a spatiotemporal selection
algorithm, refits monthly per metropolis, and is validated on \textbf{over 33 million
hourly observations} from the largest US LCS network across ten metropolitan areas. The
equity finding is the part that should be read twice: LCS spatial density varies
significantly by race/ethnicity and income, with \textbf{African American communities the
least covered} --- so a remote-calibration method that works is a prerequisite for, not a
substitute for, fixing where the sensors physically are. Note what the paper does
\emph{not} report in its abstract: no $R^2$, no RMSE, no bias figure, only
``good agreement''. A calibration paper whose headline claim is agreement and which
declines to quantify it cannot be benchmarked against the colocated literature, and this
digest will not assert a performance number it was not given.

Against that, Zhang et al.\ put a mask-format sampler on \textbf{12 volunteers} in subway
environments and collect inhaled and exhaled aerosol simultaneously. Roughly
\textbf{66\% of inhaled PM is deposited}, and deposition is highest below
\textbf{1.0\,$\mu$m}. Every ambient network in the first paper --- and every regulatory
station it is calibrated against --- reports PM$_{2.5}$ \emph{mass}, a metric to which the
sub-micron fraction contributes little while carrying most of the deposited dose. The two
papers are not in tension; they mark the two ends of a gap. We are getting much better at
making a dense network agree with a reference instrument, and no better at making that
reference quantity the one that reaches the lung.
\end{keybox}

\vspace{3mm}

% -------------------------------------------------------------- exec brief
\section*{\sffamily\bfseries\color{Deep}What changed today}
\vspace{-1.5mm}{\color{Amber}\rule{\linewidth}{1.1pt}}\vspace{2.5mm}

\footnotesize
\begin{itemize}
\item \textbf{The strongest causal-design evidence of the watch to date lands on a clean
air zone.} Wood et al.\ (\textit{Lancet Public Health}) recruited \textbf{3,414 children}
across 84 schools in London and Luton \emph{before} ULEZ implementation and followed
FEV$_1$ growth in a parallel-cohort design. London children gained
\textbf{10\,mL/year more (95\% CI 5--15, $p=0{\cdot}0002$)} than Luton children, against a
faster NO$_2$ decline ($-3{\cdot}77$ vs $-1{\cdot}77\,\mu$g/m$^3$ per year), and a baseline
London deficit of $-38$\,mL reached parity after four years. The design is the point:
a contemporaneous control city converts a policy rollout into something close to a natural
experiment. The limitation is that the exposure contrast is carried by NO$_2$, the
traffic-exhaust marker, not by PM --- the paper cannot separate the two.

\item \textbf{Two independent records argue that PM$_{2.5}$ mass is the wrong exposure
metric, from opposite directions.} Parmar et al.\ resolve Delhi PM$_{2.5}$ into components
and find \textbf{black carbon at 11.5\% (3.3--20.4) per 10\,$\mu$g/m$^3$} against
\textbf{0.93\% (0.10--1.76) for total mass} --- an order-of-magnitude difference in
per-unit potency. Li et al.\ (\textit{ES\&T}) reach the same place from the bench,
identifying chromium and specific condensed water-soluble organic constituents as the
fractions that actually cross a human blood--brain-barrier platform. Component-resolving
instrumentation is being asked for by both literatures simultaneously.

\item \textbf{A null worth carrying.} Zewdie et al.\ find wildfire smoke PM$_{2.5}$
\emph{not} associated with tract-level cancer history prevalence across the contiguous US.
The two surviving regional signals are small (Midwest \textbf{PR 1.020}, Northeast
\textbf{PR 1.040}) and sit in the \emph{least}-smoke regions --- the pattern residual
confounding produces, as the authors themselves note.

\item \textbf{A widely-embedded aerosol-software default is shown to be wrong most of the
time it is used.} Rapp et al.\ (\textit{AMT}) fit a ten-year Storm Peak size-distribution
record with a template-free lognormal decomposition and find \textbf{four or more modes
required for 55.7\% of new-particle-formation data} --- against the tri-modal assumption
standard in PSD fitting. \textit{AMT} is not PubMed-indexed and Crossref-by-ISSN carried
it, as it has carried the sensing block on most recent issues.

\item \textbf{Europe PMC contributed nothing, and it is indexing lag, not a broken leg.}
A control query returns \textbf{3 records of any topic} created on 18 August against 2,794
on 17 August --- an argument for backfilling a missed day rather than recording a null.
\end{itemize}

\clearpage

% ==================================================================== ANALYTICS
\section*{\sffamily\bfseries\color{Deep}Corpus analytics}
\vspace{-1.5mm}{\color{Amber}\rule{\linewidth}{1.1pt}}\vspace{3mm}

\noindent\includegraphics[width=\linewidth]{\FIGDIR/f1_subtopics.png}
\figcap{Subtopic assignment of the 18 carried records, single-label by dominant
contribution. \textit{Sensing, forecasting \& instrumentation} (5) and
\textit{Exposure assessment \& modelling} (4) together account for exactly half the issue;
that share has now been at or above 0.50 in seven of the last twelve issues, so the
instrumentation weighting is a standing feature of this watch rather than a property of
today. Eight of ten clusters are populated. \textit{Cardiovascular \& metabolic} and
\textit{Occupational \& indoor} are both empty --- the third issue in which that pairing is
simultaneously absent, after 4 and 16 August.}

\vspace{2mm}
\noindent\includegraphics[width=\linewidth]{\FIGDIR/f2_design_endpoint.png}
\figcap{Left --- study architecture. No single design dominates: prospective cohorts (3)
and machine-learning models (3) lead, and the tail is long, which is what a corpus split
evenly between epidemiology and instrumentation looks like. Right --- health endpoint by
organ system. Seven records carry an explicit no-health-endpoint label --- the direct
consequence of the instrumentation share noted above --- and a further two are the
abstract-withheld records, whose endpoint is unknown rather than absent.}

\vspace{2mm}
\noindent\includegraphics[width=\linewidth]{\FIGDIR/f3_metric_geography.png}
\figcap{Left --- particle metric. Note the two records whose exposure is explicitly
\emph{not} PM$_{2.5}$ mass: Zhang et al.\ resolve deposition by size below 1\,$\mu$m, and
Rapp et al.\ work in number-size-distribution space entirely. Right --- geographic
provenance. The US and China tie at six records each; South Asia contributes two (Delhi,
Bangladesh) and Latin America one, the first Argentine record of the watch --- which
required adding an \texttt{Argentina} key and needle to the geography map, without which it
would have fallen silently to the unmapped bucket.}

\vspace{2mm}
\noindent\includegraphics[width=\linewidth]{\FIGDIR/f6_lifecourse.png}
\figcap{Life-course windows addressed; counts are non-exclusive. Working-age adulthood
dominates (6), carried by the Delhi and Korean admission series, the SWAN midlife cohort
and the 12 subway volunteers. The prenatal window (3) ties 15 August for the highest count
of the last twelve issues, but the composition differs: here it is one prenatal-exposure
cohort plus a cell-line spermatocyte model, not the block of reproductive birth cohorts
that drove the mid-August run.}

\vspace{2mm}
\noindent\includegraphics[width=\linewidth]{\FIGDIR/f4_heatmap.png}
\figcap{Subtopic $\times$ study-architecture cross-tabulation. The structural feature of
this issue is visible as an empty region rather than a full one: the
\textit{Sensing} and \textit{Exposure} rows contain no design that carries a health
endpoint. Running the converse check finds a single exception --- Perera et al.\ monitored
airborne PAH personally within the CCCEH cohorts, and is the only health-endpoint record
here whose exposure was measured by the study team rather than taken from a dispersion
model, a reanalysis or a regulatory network. The wearable sampler of Zhang et al.\ closes
that gap instrumentally but carries no health outcome.}

\vspace{2mm}
\noindent\includegraphics[width=\linewidth]{\FIGDIR/f5_forest.png}
\figcap{The nine interval-bearing estimates in today's corpus on a common logarithmic
axis. \textbf{Exposure contrasts are not mutually comparable}: the Kang et al.\ and
Parmar et al.\ entries are percentage changes converted to risk-ratio equivalents (per IQR
and per 10\,$\mu$g/m$^3$ respectively), Zewdie et al.\ report prevalence ratios per IQR,
and the Jang et al.\ entries are accelerated-failure-time \emph{time ratios} for a
top-versus-bottom quartile contrast --- these last are below the null by construction and a
value under 1 means \emph{earlier} menopause, not lower risk. Read direction and precision
only. The tightest intervals belong to the two largest administrative series
(3.7 million osteoarthritis visits; all US census tracts) and the widest to the
component-resolved black-carbon estimate, which is the one with the smallest effective
exposure contrast.}

\clearpage

% ==================================================================== DIGEST
\section*{\sffamily\bfseries\color{Deep}Paper-level digest}
\vspace{-1.5mm}{\color{Amber}\rule{\linewidth}{1.1pt}}\vspace{1mm}

{\fontsize{7.4}{9}\selectfont\color{Slate}Ordered by cluster. Each entry gives the
methodological core and the finding that matters, not an abstract paraphrase.\par}

\band{Sensing, forecasting \& instrumentation\hfill 5 records}{Amber}

\paperentry{Wang et al.\ 2026}
{Remote calibration of a nationwide dense low-cost sensor network for hyperlocal air surveillance}
{Environ Sci Technol}
{Removes physical colocation from the LCS calibration loop. A spatiotemporal pairing
algorithm selects representative training data linking uncolocated sensors to reference
stations, with metropolis-specific monthly model refresh; validated on
\textbf{$>$33 million hourly observations} across ten US metros. Demonstrated downstream on
high-granularity land-use regression and wildfire plume tracking. Two things to hold
against it. First, the abstract reports agreement qualitatively and gives \emph{no}
$R^2$, RMSE or bias, so it cannot be placed against the colocated calibration literature
without the full text. Second, and more usefully, the paper quantifies its own
denominator problem: LCS density differs significantly by race/ethnicity and income, with
African American communities least covered. Remote calibration widens the usable network;
it does not move a sensor into a neighbourhood that has none.}
{10.1021/acs.est.6c03270}

\paperentry{Zhang et al.\ 2026}
{Wearable device for direct assessment of particle deposition in the human respiratory system}
{Environ Sci Technol}
{Face-mask-format sampler collecting inhaled and exhaled aerosol simultaneously, so
deposition is measured by difference rather than inferred from a lung-deposition model.
Bench validation confirmed exhaled-particle capture with minimal ambient interference;
deployed on \textbf{12 healthy volunteers} in subway environments.
\textbf{$\sim$66\% of inhaled PM deposited}, highest for \textbf{PM$_1$}; Fe, Cu and Pb
dominated the subway aerosol and deposition varied with lung function and exercise habit.
Twelve subjects in one microenvironment is a feasibility study, not a dosimetry dataset,
and the size-resolved difference method is vulnerable to hygroscopic growth of exhaled
particles at 100\% RH --- which the abstract does not address. The device concept is
nonetheless the most directly relevant instrumentation record of the issue: it measures
the delivered dose that every ambient network is a proxy for.}
{10.1021/acs.est.6c05543}

\paperentry{Rapp et al.\ 2026}
{Curve fitting algorithm for multimodal particle size distributions: a theoretical basis}
{Atmos Meas Tech}
{Open-source lognormal decomposition with no predefined mode templates and no mode-count
constraint. Each mode's $N_i$, $\sigma_g$ and $\bar{D}_{pg}$ are fitted by
Levenberg--Marquardt and iteratively subtracted until convergence. Evaluated against a
ten-year ambient record at Storm Peak Laboratory and against laboratory PSDs of
size-selected ammonium sulfate coated in oxidised biogenic SOA. The result that matters:
\textbf{a tri-modal assumption is inadequate for 55.7\% of new-particle-formation data},
four or more modes being required --- so the default configuration of most PSD fitting
software is wrong in the majority of the regime it is applied to. Agreement with existing
models at a tri-modal setting is strong with substantially lower computation. Supplemental
routines format output from common mobility particle size spectrometers, which lowers the
adoption cost considerably.}
{10.5194/amt-19-5373-2026}

\paperentry{Li et al.\ 2026}
{Crowdsourced, community-engaged VOC monitoring and source apportionment near a landfill}
{Environ Sci Technol}
{Community citizen-science campaign near the Bristol, Virginia solid-waste landfill using
a modified ultrasonic personal air sampler under EPA Method TO-17. BTEX remained below EPA
reference concentrations throughout but was significantly elevated relative to comparison
locations. Gas-phase, not particulate --- carried here for the deployment methodology
rather than the analyte. The transferable content is the study-design integration of
resident-reported odour events into sampling triggers, which is the same problem a
community PM network faces when episodic sources drive the exposure and a fixed duty cycle
misses them.}
{10.1021/acs.est.6c04657}

\paperentry{Saravagna Sai Lakshmi et al.\ 2026}
{CORTA-Net: hybrid deep learning with attention-based interpretability for PM$_{2.5}$ forecasting}
{J Vis Exp}
{CorrXGBoost-Rank feature screening, transfer-learning LSTM encoder, multi-head attention
for feature- and timestep-level attribution, over CPCB/DPCC station data, IMD meteorology
and MODIS fire counts for Delhi 2012--2024. Reported to reduce prediction error against
random forest, XGBoost, LSTM and attention-LSTM baselines --- but \textbf{no RMSE, MAE or
$R^2$ appears in the abstract for either the model or the baselines}, so the size of the
improvement is unstated and unverifiable here. A methods-video venue rather than a
modelling journal, and the contribution is an assembly of existing components. Included
for the fire-count integration, which is the one input choice that is specific to the
Delhi problem rather than generic.}
{10.3791/71004}

\band{Exposure assessment \& modelling\hfill 4 records}{Amber}

\paperentry{Dias et al.\ 2026}
{Micro-scale spatial patterns of PM$_{2.5}$ across urban green and non-green areas in an arid city}
{Bull Environ Contam Toxicol}
{Six mobile monitoring campaigns along a fixed 4.3\,km transect with
\textbf{75 georeferenced sampling points} in San Juan, Argentina, winter 2026, classified
green versus non-green by surrounding vegetation cover. Generalised linear mixed models
adjusting for traffic intensity, distance to nearest road and repeated measurement give
\textbf{$\sim$31\% lower adjusted PM$_{2.5}$ in green spaces}. The adjustment set is the
weak point --- traffic intensity and distance-to-road are the obvious confounders but
green spaces in an arid city also differ in surface dust availability and irrigation, which
is an unadjusted pathway that would produce the same sign. A single winter and a single
transect. Useful mainly as a demonstration that mobile monitoring resolves within-block
gradients that a regulatory network cannot, in a South American arid setting where almost
no such data exist.}
{10.1007/s00128-026-04318-8}

\paperentry{Jin et al.\ 2026}
{Hourly 2\,km Himawari-8 aerosol fine-mode fraction retrieval improves PM$_{2.5}$ mapping}
{Atmos Pollut Res}
{\textbf{Metadata only.} Elsevier withheld the abstract and the record is too recent for
Semantic Scholar (404) or Scholar Gateway, whose corpus currently ends in May 2026. The
registered title asserts that adding a fine-mode-fraction retrieval to Himawari-8
geostationary observations improves surface PM$_{2.5}$ estimation over the
Beijing--Tianjin--Hebei region at hourly, 2\,km resolution. \textbf{No performance figure,
validation design or comparison baseline is asserted here, because none was available.}
Flagged for retrieval on a later run. Geostationary fine-mode retrieval is the most direct
satellite analogue of the composition-resolution problem raised elsewhere in this issue,
so the numbers are worth chasing.}
{10.1016/j.apr.2026.103180}

\paperentry{Zaman \& Mim 2026}
{Decadal variation of atmospheric PM$_{2.5}$ over Bangladesh from MERRA-2 reanalysis}
{Environ Sci Atmos}
{A 25-year reanalysis climatology of PM$_{2.5}$ and its speciated components across
contrasting Bangladeshi microenvironments, presented by the authors as the first
comprehensive treatment for a country that ranks among the most PM$_{2.5}$-polluted
globally. The Crossref abstract is truncated at source after the framing sentence, so
trend magnitudes, speciation shares and the validation against ground observations are
all unavailable here. The methodological caveat is generic but load-bearing for this
kind of study: MERRA-2 assimilates AOD, not surface PM, and its skill over South Asia is
known to be poorer than over North America or Europe --- so a reanalysis-only trend for
Bangladesh needs its ground-truth comparison read before the numbers are used.}
{10.1039/d6ea00108d}

\paperentry{Guo et al.\ 2026}
{Meteorology-driven long-range transport to an ultra-low emission region during the monsoon}
{Atmos Environ}
{\textbf{Metadata only.} Elsevier withheld the abstract; Semantic Scholar returns 404 for
the DOI. The registered title indicates a Hainan case study of monsoon-driven long-range
transport into a low-background region. Nothing beyond the title is asserted.
Transport-into-clean-background studies are the useful test case for source attribution
because the local emission term is small, which is why this is carried rather than
rejected --- but it carries no extractable content until the abstract is released.}
{10.1016/j.atmosenv.2026.122293}

\band{Respiratory \& allergic\hfill 1 record}{Sage}

\paperentry{Wood et al.\ 2026}
{Ultra Low Emission Zone and lung function growth in children: prospective parallel cohort}
{Lancet Public Health}
{The strongest design in the issue. \textbf{3,414 children} recruited across 84 schools
(44 London, 40 Luton) between June 2018 and April 2019 --- \emph{before} ULEZ implementation
--- with \textbf{3,209 (94.0\%) in the FEV$_1$ analysis}. Luton functions as a
contemporaneous control city, converting the policy rollout into a difference-in-differences
natural experiment rather than a before--after series. Baseline London deficit
\textbf{$-38$\,mL (95\% CI $-58$ to $-18$)} against a baseline NO$_2$ excess of
\textbf{18.95\,$\mu$g/m$^3$}; over follow-up London FEV$_1$ growth exceeded Luton by
\textbf{10\,mL/year (5--15, $p=0{\cdot}0002$)} (233 vs 223\,mL/year) and mean FEV$_1$
reached parity at four years. Clinically impaired lung function fell 14\%$\to$9\% in
London and 9\%$\to$7\% in Luton. Two constraints on reading it as a PM result: the
exposure contrast is expressed in NO$_2$, the exhaust marker, so PM and NO$_2$ effects are
not separable; and two cities differ in more than their air, which parallel-cohort
adjustment mitigates but does not eliminate.}
{10.1016/S2468-2667(26)00163-5}

\band{Neuro, cognition \& mental health\hfill 2 records}{Deep}

\paperentry{Parmar et al.\ 2026}
{Short-term PM$_{2.5}$ and its chemical components and neuropsychiatric hospitalisations}
{Environ Int}
{Ground observations fused with 5\,km WRF-Chem output to build a gridded PM$_{2.5}$ field
for Delhi 2016--2019, then GAM with distributed lags against hospital neuropsychiatric
admissions. Total mass gives \textbf{+0.926\% (95\% CI 0.097--1.762) per 10\,$\mu$g/m$^3$};
\textbf{black carbon gives +11.509\% (3.272--20.402)}, ahead of organic carbon and other
inorganics, and only industrial-source PM$_{2.5}$ was significant among source
contributions. Among disorder categories only anxiety reached significance. The
per-unit-mass potency gap between BC and total mass is the headline, but it should be read
with the exposure ranges in mind: BC concentrations are roughly an order of magnitude
smaller than total mass in Delhi, so a per-10\,$\mu$g/m$^3$ contrast for BC is a large
extrapolation beyond the observed spread. Single-hospital case ascertainment also limits
generalisation across a city of Delhi's size.}
{10.1016/j.envint.2026.110464}

\paperentry{Perera et al.\ 2026}
{Declining prenatal air pollution and improving cognitive scores across successive NYC cohorts}
{Environ Health}
{Three successive Columbia CCCEH birth cohorts, 1998--2020, \textbf{n = 1,109}, with
prenatal PM$_{2.5}$/NO$_2$ (first two cohorts) and airborne PAH (all three), and Bayley
MDI at ages 1, 2 and 3. Exposures fell significantly across cohorts and MDI rose
significantly; pooled associations were inverse and significant at all three ages, and
cohort-specific inverse associations were clearest in the earliest, highest-exposure
cohort ($n > 700$) and weaker in the 2008--2016 cohort. The design is the interesting
part --- using secular exposure decline as the contrast --- and also the vulnerability:
everything else about being a child in New York changed between 1998 and 2020, and
period effects on Bayley scores are not identifiable from cohort effects here. The
authors frame it as policy-relevant evidence; it is better read as consistency between two
trends than as an estimated effect.}
{10.1186/s12940-026-01328-1}

\band{Reproductive \& developmental\hfill 1 record}{Sky}

\paperentry{Jang et al.\ 2026}
{Ambient air pollution and age at natural menopause in a multi-ethnic US cohort}
{J Clin Endocrinol Metab}
{SWAN, \textbf{1,457 midlife women} across six US sites 1999--2015, five racial/ethnic
groups, with 1\,km$^2$ annual PM$_{2.5}$, NO$_2$ and O$_3$ linked to residence and age at
natural menopause modelled as time-to-event by accelerated failure time.
\textbf{NO$_2$ TR 0.983 (0.971--0.995), O$_3$ TR 0.983 (0.971--0.996)} for the highest
versus lowest exposure quartile --- roughly one year earlier menopause, which the authors
benchmark against the 0.8-year advance seen in current smokers.
\textbf{PM$_{2.5}$ was not significantly associated}, which is the finding this watch
should record plainly rather than fold into a general air-pollution claim. Residential
address as the exposure surrogate over a 16-year window in a mobile midlife population is
the main measurement limitation.}
{10.1210/clinem/dgag340}

\band{Mechanistic toxicology\hfill 2 records}{Violet}

\paperentry{Li et al.\ 2026b}
{Component-resolved PM$_{2.5}$ translocation across a human blood--brain-barrier immune platform}
{Environ Sci Technol}
{Microfluidic BBB--immune microenvironment platform with compartment-resolved
quantification of which PM$_{2.5}$ constituents actually cross, rather than inference from
downstream effect. High doses increased barrier permeability and drove constituent
translocation; translocated material produced astrocyte reactivity, microglial M1
polarisation, elevated ROS, neuroinflammation and lipid-metabolic dysregulation.
Component-resolved correlation implicates \textbf{chromium} and the water-soluble organic
matter fraction --- specifically the condensed species C$_{28}$H$_7$O$_2$N and
C$_{21}$H$_{16}$N$_2$SP$_2$. The dose is described only as ``high'', with no
concentration given in the abstract and therefore no way to judge ambient relevance, which
is the standard failure mode of this literature. The compartment-resolved measurement of
translocation is nonetheless a genuine methodological advance over effect-only designs.}
{10.1021/acs.est.6c01030}

\paperentry{Jin \& Zhang 2026}
{Astragaloside IV protects against PM$_{2.5}$-induced spermatocyte injury via cGAS-STING/Nrf2/HMOX1}
{Reprod Toxicol}
{Network pharmacology and GEO reanalysis (GSE189187) to nominate targets, molecular
docking, then validation in GC-2 spermatocytes. PM$_{2.5}$ at
\textbf{200\,$\mu$g/mL} raised ROS 5.35-fold and apoptosis to 42.83\%; AS-IV co-treatment
restored viability to 83.3\%, suppressed cGAS/p-TBK1/p-IRF3 and activated Nrf2/HMOX1,
reducing apoptosis to 18.27\%. The docking affinities exceed the reference inhibitors
G150 and H-151, which is presented as validation and is not --- docking scores do not
rank cellular potency. The exposure is a cell-culture suspension four orders of magnitude
above any inhaled concentration, and \textbf{AUC = 1.000} for a single gene in a public
transcriptomic dataset is a small-sample artefact rather than a discriminative result.
Carried at tier C for the pathway hypothesis only.}
{10.1016/j.reprotox.2026.109328}

\band{Burden, policy \& mitigation\hfill 1 record}{Coral}

\paperentry{Liu et al.\ 2026}
{Disproportionate effects of direct air capture on the regional health cobenefits of net-zero}
{Environ Sci Technol}
{Grid-level (36\,km$^2$) air-pollution mortality assessment along China's net-zero CO$_2$
pathway 2020--2060, contrasting scenarios with and without large-scale DACCS deployment.
DACCS substitutes for fossil-fuel phase-out and therefore forgoes the co-pollutant
reduction that would otherwise occur: \textbf{1.75 million additional cumulative premature
deaths}, with the losses concentrated in developing provinces
(\textbf{Henan $-$0.13 million and Hebei $-$0.12 million avoided deaths}) via more active
emissions trading. The paradox the authors report honestly is that inter-regional
\emph{mortality-rate} inequality narrows while the distribution of forgone benefit widens.
Scenario-model output over a 40-year horizon, so the intervals are structural-uncertainty
bounded rather than statistical, and the result is a property of the assumed trading
architecture as much as of the technology.}
{10.1021/acs.est.6c04728}

\band{Other clinical endpoints\hfill 2 records}{Slate}

\paperentry{Kang et al.\ 2026}
{Short-term ambient air pollution and acute osteoarthritis morbidity across subgroups}
{Environ Pollut}
{Time-stratified case-crossover over \textbf{3,697,314 OA-related hospital visits}
(ICD-10 M15--M19) from the Korean NHIS-NSC 2002--2019, linked to national pollutant fields,
with health-screening data enabling individual-level effect modification. Per IQR:
\textbf{NO$_2$ +5.28\% (5.02--5.54)}, CO +2.88\%, SO$_2$ +1.36\%,
\textbf{PM$_{10}$ +0.90\% (0.73--1.07)}, O$_3$ +0.76\%,
\textbf{PM$_{2.5}$ +0.73\% (0.47--1.00)} --- so the particulate terms are the weakest of the
six and NO$_2$ is the only one robust in two-pollutant models. Enormous power makes
intervals tight enough to exclude the null on effects this small, which is exactly when
outcome-definition error matters: the authors state plainly that claims-based visit dates
index healthcare utilisation, not symptom onset. A same-day traffic-pollution association
with a discretionary outpatient visit is at least as consistent with behaviour as with
pathophysiology.}
{10.1016/j.envpol.2026.128985}

\paperentry{Zewdie et al.\ 2026}
{Wildfire smoke PM$_{2.5}$ and cancer history prevalence at census-tract level}
{Cancer Epidemiol}
{Daily wildfire-specific smoke PM$_{2.5}$ aggregated to cumulative (2006--2020) and recent
(2018--2020) means, linked to 2022 tract-level cancer history prevalence for all
contiguous-US tracts; quasi-Poisson GAMs with random effects and spatial splines, adjusted
for sociodemographics, behaviours, health status, density, climate and region.
\textbf{No association in fully adjusted models for either exposure window.} Regional
effect modification produced small positive estimates in the Midwest
(\textbf{PR 1.020, 1.013--1.026}) and Northeast (\textbf{PR 1.040, 1.032--1.049}) --- the
two regions with the least smoke exposure, a spatial pattern more consistent with residual
confounding than with dose. Cross-sectional and ecological, with prevalence of \emph{ever}
having cancer set against a 15-year exposure window that cannot precede diagnosis for most
cases; latency alone makes this design close to uninformative for incidence. The value
here is a well-powered, honestly reported null in a literature prone to amplification.}
{10.1016/j.canep.2026.103199}

\clearpage

% ==================================================================== TABLE
\section*{\sffamily\bfseries\color{Deep}Corpus register}
\vspace{-1.5mm}{\color{Amber}\rule{\linewidth}{1.1pt}}\vspace{2.5mm}

{\fontsize{6.7}{8.4}\selectfont
\setlength{\tabcolsep}{2.4pt}
\begin{longtable}{@{}>{\RaggedRight}p{25mm}>{\RaggedRight}p{41mm}>{\RaggedRight}p{27mm}%
>{\RaggedRight}p{23mm}>{\RaggedRight}p{18mm}>{\RaggedRight}p{15mm}@{}}
\toprule
\textsf{\textbf{First author}} & \textsf{\textbf{Focus}} & \textsf{\textbf{Journal}} &
\textsf{\textbf{Design}} & \textsf{\textbf{Metric}} & \textsf{\textbf{Region}}\\
\midrule
\endfirsthead
\toprule
\textsf{\textbf{First author}} & \textsf{\textbf{Focus}} & \textsf{\textbf{Journal}} &
\textsf{\textbf{Design}} & \textsf{\textbf{Metric}} & \textsf{\textbf{Region}}\\
\midrule
\endhead
\multicolumn{6}{@{}l@{}}{\cellcolor{Mist}\textsf{\textbf{\textcolor{Deep}{Neuro / mental health}}}}\\[0.6mm]
\href{https://doi.org/10.1001/jamanetworkopen.2026.25015}{Dales} & Wildfire-sourced PM$_{2.5}$ and ED visits for migraine and other primary headache & \textit{JAMA Netw Open} & Case-crossover & PM$_{2.5}$ (wildfire) & Canada\\
\href{https://doi.org/10.1093/ageing/afag215}{Glover} & Umbrella meta-review of urban environmental exposures and cognitive health & \textit{Age Ageing} & Umbrella review / meta-analysis & PM$_{2.5}$ & Global\\
\href{https://doi.org/10.1016/j.nsa.2026.107019}{Goossens} & In utero air pollution exposure and pre-/postnatal brain development & \textit{Neurosci Appl} & Narrative review & PM$_{2.5}$ / PM$_{10}$ / PM$_{\mathrm{coarse}}$ & Global\\
\href{https://doi.org/10.1016/j.isci.2026.116617}{Ma} & PM$_{2.5}$-induced neurotoxicity and depression: exposure risks and mechanisms & \textit{iScience} & Narrative review & PM$_{2.5}$ + components & Global\\
\href{https://doi.org/10.3760/cma.j.cn112150-20250804-00752}{Pan} & Maternal PM exposure and offspring neurodevelopment: effects and mechanisms & \textit{Zhonghua Yu Fang Yi Xue Za Zhi} & Narrative review & PM (bulk) & China\\
\addlinespace[0.9mm]
\multicolumn{6}{@{}l@{}}{\cellcolor{Mist}\textsf{\textbf{\textcolor{Teal}{Cardiovascular \& metabolic}}}}\\[0.6mm]
\href{https://doi.org/10.3760/cma.j.cn112338-20251120-00843}{Huang} & PM$_{2.5}$ and dyslipidaemia in rural Chinese children; dietary effect modification & \textit{Zhonghua Liu Xing Bing Xue Za Zhi} & Cross-sectional (multistage) & PM$_{2.5}$ & China\\
\href{https://doi.org/10.32394/pe/218966}{Kazmierska} & Air pollution and cardiovascular disease: review of current data & \textit{Przegl Epidemiol} & Narrative review & PM$_{2.5}$ / PM$_{10}$ & Global\\
\href{https://doi.org/10.1007/s00125-026-06814-2}{Xi} & Air pollution, built environment and temperature vs. prediabetes phenotypes & \textit{Diabetologia} & Prospective cohort & PM$_{2.5}$ & Germany\\
\href{https://doi.org/10.1161/JAHA.125.046917}{Yang} & Short-term PM$_{2.5}$ and hospitalisation/ED use in atrial fibrillation & \textit{J Am Heart Assoc} & Case-crossover & PM$_{2.5}$ (wildfire-influenced) & USA\\
\addlinespace[0.9mm]
\multicolumn{6}{@{}l@{}}{\cellcolor{Mist}\textsf{\textbf{\textcolor{Sky}{Reproductive \& developmental}}}}\\[0.6mm]
\href{https://doi.org/10.1016/j.jes.2025.10.043}{Chen} & Prenatal air pollution, fetal growth trajectory phenotypes and child neurodevelopment & \textit{J Environ Sci (China)} & Prospective cohort & PM$_{2.5}$ / PM$_{10}$ & China\\
\href{https://doi.org/10.1016/j.ecoenv.2026.120566}{Huang} & Joint prenatal bisphenol + air-pollutant mixture exposure and preschool IQ & \textit{Ecotoxicol Environ Saf} & Prospective cohort (mixtures) & PM$_{2.5}$ / PM$_{10}$ & China\\
\href{https://doi.org/10.1007/s40572-026-00552-8}{Peterson} & Ambient air pollution, PAHs and pubertal development: critical review & \textit{Curr Environ Health Rep} & Critical review & PM$_{2.5}$ components & Global\\
\href{https://doi.org/10.1016/j.envres.2026.125313}{Xie} & Long-term PM$_{2.5}$ and timing of pubertal development in boys and girls & \textit{Environ Res} & Prospective cohort & PM$_{2.5}$ & China\\
\addlinespace[0.9mm]
\multicolumn{6}{@{}l@{}}{\cellcolor{Mist}\textsf{\textbf{\textcolor{Sage}{Respiratory \& allergic}}}}\\[0.6mm]
\href{https://doi.org/10.1007/s00484-026-03289-z}{Bayar Muluk} & Climate change and otolaryngology: airway disease impacts and emission reduction & \textit{Int J Biometeorol} & Narrative review & PM$_{2.5}$ & Global\\
\href{https://doi.org/10.4103/mgr.MEDGASRES-D-25-00164}{Liao} & Short-term air pollution, respiratory symptoms and lung function in asthma & \textit{Med Gas Res} & Panel study & PM$_{2.5}$ / PM$_{10}$ & China\\
\href{https://doi.org/10.3390/epidemiologia7040097}{Sanchez de Prada} & Ambient PM and asthma-coded admissions in children 0-3 y, Spain & \textit{Epidemiologia} & Ecological panel & PM$_{2.5}$ / PM$_{10}$ & Spain\\
\href{https://doi.org/10.1016/j.jhazmat.2026.143076}{Yang} & Allergenic vs. toxic risk of airborne seed hairs compared with ambient PM$_{2.5}$ & \textit{J Hazard Mater} & Field measurement & PM$_{2.5}$ + bioaerosol & China\\
\addlinespace[0.9mm]
\multicolumn{6}{@{}l@{}}{\cellcolor{Mist}\textsf{\textbf{\textcolor{Violet}{Mechanistic toxicology}}}}\\[0.6mm]
\href{https://doi.org/10.1016/j.jes.2025.12.074}{Ding} & PM$_{2.5}$ promotes vascular calcification via miR-27b-5p/HES1 signalling & \textit{J Environ Sci (China)} & In vitro + organoid + murine & PM$_{2.5}$ & China\\
\href{https://doi.org/10.1007/s00360-026-01698-5}{Maraschi} & Osmoregulatory disruption and metal bioaccumulation from settleable atmospheric PM & \textit{J Comp Physiol B} & Controlled exposure (ecotox) & Settleable APM & Brazil\\
\href{https://doi.org/10.1016/j.jes.2025.11.018}{Wu} & Macrophage NLRP3 inflammasome in real-ambient PM-induced hepatic glucose dysregulation & \textit{J Environ Sci (China)} & Real-ambient murine + co-culture & Real-ambient PM & China\\
\addlinespace[0.9mm]
\multicolumn{6}{@{}l@{}}{\cellcolor{Mist}\textsf{\textbf{\textcolor{Clay}{Occupational \& indoor}}}}\\[0.6mm]
\href{https://doi.org/10.3389/fmedt.2026.1796828}{Chen} & Nested suction instrument to cut intraoperative surgical-smoke PM$_{2.5}$ & \textit{Front Med Technol} & Bench / device experiment & PM$_{2.5}$ (surgical smoke) & China\\
\href{https://doi.org/10.1093/annweh/wxag055}{Cole} & Determinants of respirable crystalline silica exposure in Australian tunnelling & \textit{Ann Work Expo Health} & Observational exposure study & Respirable dust / RCS & Australia\\
\href{https://doi.org/10.1016/j.jhazmat.2026.143081}{Ding} & Stage-resolved hair Zn-isotope dynamics as a firefighting combustion-exposure tracer & \textit{J Hazard Mater} & Repeated-measures biomarker & Combustion-derived PM & China\\
\href{https://doi.org/10.3390/jox16040119}{Ramsden} & Xenobiotic hazards in aircraft cabin air & \textit{J Xenobiot} & Review + risk economics & Ultrafine particles & Global\\
\addlinespace[0.9mm]
\multicolumn{6}{@{}l@{}}{\cellcolor{Mist}\textsf{\textbf{\textcolor{Amber}{Sensing, forecasting \& instrumentation}}}}\\[0.6mm]
\href{https://doi.org/10.1016/j.jenvman.2026.130515}{Shi} & Adaptive spatiotemporal graph transformer for multi-site multi-step PM$_{2.5}$ forecasting & \textit{J Environ Manage} & Deep-learning model & PM$_{2.5}$ & China / India\\
\href{https://doi.org/10.3390/bios16070380}{Tsuchiya} & Seasonal and indoor-outdoor characteristics of fluorescent aerosol particles in offices & \textit{Biosensors} & Real-time sensor campaign & Bioaerosol / coarse & Japan\\
\addlinespace[0.9mm]
\multicolumn{6}{@{}l@{}}{\cellcolor{Mist}\textsf{\textbf{\textcolor{Amber}{Exposure assessment \& modelling}}}}\\[0.6mm]
\href{https://doi.org/10.1029/2025GH001621}{Solorzano-Araque} & Multiscale exposure modelling: LOTOS-EUROS + mobility surveys + particle toxicity & \textit{GeoHealth} & CTM + mobility integration & PM$_{2.5}$ / PM$_{10}$ & Colombia\\
\addlinespace[0.9mm]
\multicolumn{6}{@{}l@{}}{\cellcolor{Mist}\textsf{\textbf{\textcolor{Coral}{Burden, policy \& mitigation}}}}\\[0.6mm]
\href{https://doi.org/10.1002/hsr2.72886}{Karamzad} & Ischaemic heart disease burden in MENA 1990-2021; PM attributable fraction 28.2\% & \textit{Health Sci Rep} & GBD secondary analysis & Ambient PM & MENA\\
\href{https://doi.org/10.1371/journal.pone.0347989}{Rees} & HEPA filtration to reduce respiratory infection in care homes & \textit{PLoS One} & RCT process evaluation & Filtered PM & UK\\
\href{https://doi.org/10.1016/j.jes.2026.01.009}{Shan} & Policy-driven emission reductions and post-pandemic rebound, high-resolution inventories & \textit{J Environ Sci (China)} & Emission inventory & PM emissions & China\\
\addlinespace[0.9mm]
\multicolumn{6}{@{}l@{}}{\cellcolor{Mist}\textsf{\textbf{\textcolor{Slate}{Other clinical endpoints}}}}\\[0.6mm]
\href{https://doi.org/10.3390/dj14070440}{Lupita} & PM$_{2.5}$, preventive-care capacity and caries/edentulism burden across the EU & \textit{Dent J} & Ecological panel & PM$_{2.5}$ & EU-27\\
\href{https://doi.org/10.1016/j.jdin.2026.06.008}{Park} & Long-term PM exposure and risk of androgenetic alopecia & \textit{JAAD Int} & Nationwide cohort & PM$_{2.5}$ / PM$_{10}$ & South Korea\\
\href{https://doi.org/10.3171/2026.3.PEDS2615}{Sun} & Environmental exposure and area deprivation vs. paediatric hydrocephalus outcomes & \textit{J Neurosurg Pediatr} & Registry cohort + geospatial & PM$_{2.5}$ (satellite-derived) & USA\\
\addlinespace[0.9mm]
\bottomrule
\end{longtable}}


\vspace{4mm}
\begin{keybox}
{\sffamily\bfseries\footnotesize\color{Deep}Provenance and method}\\[1.2mm]
{\fontsize{7.2}{8.8}\selectfont
Entry window \textbf{18 August 2026}, contiguous with the 17 August issue. This issue is a
\textbf{backfill}: the scheduled 18 August run did not execute, and the window was
recovered day-by-day on 19 August before the current day was opened.
Sources queried: \textbf{PubMed} E-utilities and the PubMed connector on two separate axes
(health, 23 records; sensing, 1, already in the health set);
\textbf{Crossref by journal ISSN} across 18 aerosol and atmospheric titles (7 deposits),
which carried the entire instrumentation block --- \textit{Atmos.\ Meas.\ Tech.},
\textit{Atmos.\ Chem.\ Phys.}\ and \textit{Atmos.\ Pollut.\ Res.}\ are not PubMed-indexed;
\textbf{Europe PMC} on \texttt{CREATION\_DATE} (0 records --- a control query returns
3 records of \emph{any} topic for 18 August against 2{,}794 for 17 August, so this is
index lag, not harvester failure); \textbf{arXiv} under a 3-day recency screen
(0 new; newest in-scope posting was 15 August and already carried);
\textbf{OpenAlex} HTTP 429, sixteenth consecutive failure, no longer retried;
\textbf{Consensus} returned three 2026 calibration papers, none deposited in-window
(ninth consecutive zero-carry, retained as a weekly/monthly gap-filler);
\textbf{ClinicalTrials.gov} 0 registrations or updates in window, \texttt{trials.json}
unchanged.
\textbf{18 records carried, 12 rejected} --- rejections were aerosol- or
geochemistry-only records with no PM monitoring method and no health endpoint (marine
organosulfate formation, iodic-acid nucleation, pyrogenic-carbon radical chemistry,
microplastics in snow, brown-carbon and siloxane SOA precursor chemistry, PTR-ToF-MS
gas-phase VOC interference, a GHG mapping framework, a vehicle-fleet material-flow model,
an OECD renewable-energy panel regression, tropical cloud Doppler measurements and a
precipitation trace-metal series). Each is logged with a reason in
\texttt{state/rejected.jsonl}.
Two records ship as \textbf{Metadata only (no abstract)}: Elsevier withheld both, Semantic
Scholar returned 404, and Scholar Gateway's corpus ends in May 2026. Nothing beyond the
registered title is asserted for either.
Subtopic, design, particle-metric and geography labels are single-label by dominant
contribution and are a judgement, not a controlled vocabulary. Effect estimates are
transcribed from abstracts and use heterogeneous exposure contrasts --- see the forest-plot
caption. All DOIs verified against Crossref. Abstracts remain the copyright of their
respective publishers.\par}
\end{keybox}

\end{document}
